\chapter{Assignment 2: Ribosome}

Ribosomes are the cellular components responsible for protein synthesis. 
Out of other side quests, the ribosome's main job is to translate the genetic information received from the mRNA into proteins (one example that comes in mind is the enzyme). 

Even though their function is the same across all living organisms, their structure is little bit different when comparing prokaryotes and eukaryotes cells. 

In bacteria (which are prokaryotes), ribosomes are smaller and are classified as 70S, being made up of a 50S and a 30S subunit. \cite{alberts2014molecular} Those cells have a nucleoid where the DNA sits.
In eukaryotic cells, however, ribosomes are larger and slightly more complex, known as 80S (consisting of 60S and 40S subunits). \cite{alberts2014molecular} These cells come with a nucleus, protected by a membrane, where the DNA is stored.

In practice, this structural difference is actually very important. For example, most of the antibiotics work by specifically targeting and binding only to the bacterial 70S ribosomes, blocking them from producing proteins. 
This stops the bacterial infection without harming the 80S ribosomes of the human cells.

The ribosome reads the mRNA sequence almost like reading lines of code. 
It processes this genetic information in small blocks of three letters at a time (those codons can be composed by the combination of the following letters: A (Adenine), U (Uracil), G (Guanine), C (Cytosine); the difference between DNA and RNA is that DNA has T (Thymine) instead of U (Uracil)). For each specific block it reads, a different helper molecule called transfer RNA (tRNA) brings a matching amino acid to the ribosome.  
The ribosome then acts like a biological compiler, linking these amino acids together one by one to form a long chain. 
It keeps running this loop until it hits a specific 'stop' codon in the sequence. 

At that exact point, the process terminates, and the newly built chain is released to fold into a functional protein." \cite{nelson2017lehninger}

In eukaryotic cells, the ribosomes can be found in two main places. 
Some of them are floating freely in the cytoplasm, which is the fluid filling the cell. These free ribosomes mostly manufacture proteins that will be used locally, right inside that same cell. 
On the other hand, many ribosomes are attached to a large membrane structure, the Endoplasmic Reticulum.
The ribosomes working on the ER usually compile proteins that are meant to be exported outside the cell or embedded into the cell membrane. \cite{steitz2008structural}

In conclusion, the ribosome is much more than just a simple cellular component; 
it acts as the essential hardware that executes the genetic code. 
By constantly reading mRNA instructions and compiling them into functional proteins.